\clearpage\fancyhead[L]{\small GAUSS TO BODE.}\hypertarget{native-404-part-lower}{}\label{native-404-part-lower}
\begin{center}\large GAUSS TO BODE.\end{center}
\begin{center}\hrule\vspace{4pt}\fontsize{10}{12}\selectfont
Astronomisches Jahrbuch für 1820. pp. 201..205. Berlin 1817.\par\vspace{4pt}\hrule\vspace{4pt}\end{center}
\begin{flushright}Göttingen, 17 March 1817.\end{flushright}
\par\indent Since the theory of the motion of Pallas is sufficiently developed, I asked Dr. \textsc{Tittel} of Erlau, who has spent a year and a half here, to calculate the ephemeris for 1818..1819 \SourceLeaf{405} now. I am pleased to enclose a copy for your next yearbook.

\par\indent Was your sky last winter as unfavorable for observations as it was here? It was very rare to catch sight of the Sun or a star. It snowed all day on 19 November last year, and nothing could be seen of the solar eclipse. Some dimming of daylight could be detected around the middle of the eclipse, but there was nothing extraordinary about it. At any rate, it has often been just as dark since then, even at noon, in snowstorms without any solar eclipse. At the lunar eclipse of 4 December the sky was somewhat less unfavorable. The beginning and the end remained invisible, but I observed the following phases quite well, in some cases between swiftly passing clouds.

\Needspace{15\baselineskip}\par\indent 4 December 1816 at the new observatory.

\par\indent Ingresses.
\begin{center}\fontsize{10}{12}\selectfont\setlength{\tabcolsep}{3pt}\begin{tabular}{llrrrl}
Tycho, first edge & $\cdot\;\cdot\;\cdot\;\cdot\;\cdot\;\cdot\;\cdot\;\cdot\;\cdot\;\cdot\;\cdot\;\cdot$ &  $8^{\mathrm u}$ & $17^{\mathrm m}$ & $2^{\mathrm s}$ & \shortstack{True\\time}\\
Tycho, second edge & $\cdot\;\cdot\;\cdot\;\cdot\;\cdot\;\cdot\;\cdot\;\cdot\;\cdot\;\cdot\;\cdot\;\cdot$ &  & $18$ & $44$ & —\\
Kepler, center & $\cdot\;\cdot\;\cdot\;\cdot\;\cdot\;\cdot\;\cdot\;\cdot\;\cdot\;\cdot\;\cdot\;\cdot\;\cdot$ &  & $29$ & $24$ & —\\
Copernicus, first edge & $\cdot\;\cdot\;\cdot\;\cdot\;\cdot\;\cdot\;\cdot\;\cdot\;\cdot\;\cdot\;\cdot$ &  & $40$ & $28$ & —\\
Copernicus, second edge & $\cdot\;\cdot\;\cdot\;\cdot\;\cdot\;\cdot\;\cdot\;\cdot\;\cdot\;\cdot\;\cdot$ &  & $45$ & $11$ & —\\
Aristarch, first edge & $\cdot\;\cdot\;\cdot\;\cdot\;\cdot\;\cdot\;\cdot\;\cdot\;\cdot\;\cdot\;\cdot$ &  & $50$ & $28$ & —\\
Dionysius, center & $\cdot\;\cdot\;\cdot\;\cdot\;\cdot\;\cdot\;\cdot\;\cdot\;\cdot\;\cdot\;\cdot\;\cdot$ &  $9$ & $0$ & $46$ & —\\
Censorinus, center & $\cdot\;\cdot\;\cdot\;\cdot\;\cdot\;\cdot\;\cdot\;\cdot\;\cdot\;\cdot\;\cdot$ &  & $5$ & $44$ & —\\
\end{tabular}\end{center}

\Needspace{7\baselineskip}\begin{flushright}Göttingen, 15 August 1817.\end{flushright}

\par\indent The construction of our new observatory is still unfinished, but it has progressed far enough for observations to be made in two rooms. Until the new fixed meridian instruments arrive and are installed, the most important observations will be those made with the 12-inch \textsc{Reichenbach} repeating circle. I am sending you the results obtained with it for Polaris and for the Sun at the last solstice. First I must note that according to my geodetic measurements the center of the new observatory lies $6''1$ farther south and $1^{\mathrm s}90$ of time farther east than the center of the old observatory. I must also note that I made a small alteration to the instrument itself. In my earlier observations I always used the brass ring fitted over the objective end of the telescope, which counterbalances both the excess weight of the prismatic eyepiece and the clamp. It had been suspected that this ring might cause the telescope to flex (or, more precisely, cause the objective end to flex more than the eyepiece causes the other end to flex). In that case the instrument would have to give all zenith distances as too \emph{small}, and it was conjectured that this might explain the difference found in determining the altitude of the pole from observations of circumpolar stars or from observations of the Sun. If that supposition were well founded, removing the counterweight would have to eliminate the difference, or rather reverse it, since the eyepiece end of the telescope now decidedly outweighs the other end. But my new observations, made with that weight removed and a weight screwed onto the alidade circle itself instead, run entirely counter to this and seem to me to prove conclusively that the flexure of the telescope is quite imperceptible. For one thing, the new observations of Polaris and the Sun give exactly the same difference in the \SourceLeaf{406}altitude of the pole and in the same direction as the earlier ones. For another, the absolute determination of the altitude of the pole at the new observatory is exactly the same as the value transferred from the old observatory.

\Needspace{15\baselineskip}\par\indent\emph{Observations of Polaris at lower culmination.}

\begin{center}\fontsize{10}{12}\selectfont\setlength{\tabcolsep}{3pt}
\begin{tabular}{l|c|c|c}
1817 & Count & \shortstack{True zenith distance\\on the meridian} & Altitude of the pole \\
\hline
March 28 & 16 & $40^\circ\ 8'\ 13''18$ & $51^\circ\ 31'\ 49''33$ \\
April\ \ 1 & 22 & $15.86$ & $47.84$ \\
\ \ \ \ \ \ \ \ 2 & 24 & $13.87$ & $50.13$ \\
\ \ \ \ \ \ \ \ 8 & 28 & $15.76$ & $50.03$ \\
\ \ \ \ \ \ \ 11 & 26 & $16.16$ & $50.51$ \\
\ \ \ \ \ \ \ 13 & 12 & $20.51$ & $46.75$ \\
May\ \ \ 7 & 30 & $23.90$ & $49.60$ \\
Mean & 158 & $\cdot\ \cdot\ \cdot\ \cdot\ \cdot$ & $51^\circ\ 31'\ 49''42$ \\
\end{tabular}
\end{center}

\par\indent In calculating refraction \textsc{Carlini}’s table was used, and for the declination of Polaris, \textsc{Bessel}’s. The observing site is $0''12$ north of the center, so the altitude of the pole at the center, from the above observations, is $51^\circ\ 31'\ 49''3$.

\par\indent My 1813 observations of Polaris at lower culmination, reduced with the same elements, give $51^\circ\ 31'\ 55''48$ for the altitude of the pole at the old observatory. The transfer thus gives the same result as the direct observations.

\par\indent The weather favored observations of the summer solstice this year more than ever. The observing site is $0''16$ south of the center.

\Needspace{23\baselineskip}\par\indent\emph{Observations of the solstice.}

\begin{center}\fontsize{10}{12}\selectfont\setlength{\tabcolsep}{3pt}
\begin{tabular}{l|c|c|c}
1817 & Count & \shortstack{True zenith distance\\on the meridian} & At the solstice \\
\hline
June\ 13 & 10 & $28^\circ\ 18'\ 11''30$ & $28^\circ\ 3'\ 52''26$ \\
\ \ \ \ \ \ 16 & 16 & $9\ \ 51.80$ & $48.63$ \\
\ \ \ \ \ \ 17 & 20 & $7\ \ 57.26$ & $55.30$ \\
\ \ \ \ \ \ 18 & 16 & $6\ \ 19.03$ & $53.46$ \\
\ \ \ \ \ \ 19 & 12 & $5\ \ \ 8.69$ & $54.66$ \\
\ \ \ \ \ \ 20 & 20 & $4\ \ 18.74$ & $51.41$ \\
\ \ \ \ \ \ 21 & 12 & $3\ \ 54.41$ & $48.91$ \\
\ \ \ \ \ \ 22 & 10 & $4\ \ \ 0.64$ & $52.10$ \\
\ \ \ \ \ \ 23 & 12 & $4\ \ 28.17$ & $51.77$ \\
\ \ \ \ \ \ 24 & 16 & $5\ \ 20.62$ & $51.62$ \\
\ \ \ \ \ \ 25 & 10 & $6\ \ 36.39$ & $49.96$ \\
\ \ \ \ \ \ 26 & 14 & $8\ \ 20.16$ & $51.61$ \\
\ \ \ \ \ \ 27 & 22 & $10\ 28.75$ & $53.49$ \\
\ \ \ \ \ \ 29 & 16 & $15\ 54.34$ & $51.86$ \\
\ \ \ \ \ \ 30 & 8 & $19\ 16.59$ & $53.88$ \\
\multicolumn{3}{l}{Mean of 214 observations} & $28^\circ\ 3'\ 52''15$ \\
\end{tabular}
\end{center}

\par\indent The zenith distance in the first column is merely reduced to the meridian and corrected for refraction. In the last column the effects of parallax and the Sun’s latitude have been removed and the values reduced to the instant of the solstice. For the center of the observatory I thus obtain $28^\circ\ 3'\ 52''31$. With the obliquity of the ecliptic $23^\circ\ 27'\ 52''9$ from the Milan ephemerides, this gives $51^\circ\ 31'\ 45''21$ for the altitude of the pole, $4''1$ less than the result from Polaris. Or, conversely, if the altitude of the pole from Polaris is taken as the basis, the obliquity is $4''1$ greater than according to \textsc{Carlini}. According to \textsc{Bessel}, the difference would be another second and a half greater. Judging from all that has been said so far about this \SourceLeaf{407}remarkable phenomenon, we do not seem to have come any closer to explaining it. It would certainly be desirable for an experienced observer to make similar observations of the Sun and of southern circumpolar stars with \textsc{Reichenbach} repeating circles at a latitude that is not too low in the \emph{southern} hemisphere. For some considerable time governments have in fact done nothing at all to support \emph{journeys} to distant parts of the world for astronomical purposes, which formerly provided us with so much valuable experience and insight, however much these same governments otherwise support science in Europe.
